\section{Supported sheaves and isolated contributions}
\label{subsec:x9-rank-one-index}

The rigid surfaces $S_-=\mathbb F_1$ and $S_+=\mathbb P^2$ also give
a computable component of a protected count. Compactify the spatial
direction of the wrapped-M5 string on a circle. The resulting type-IIA
charges are D4--D2--D0 charges. We use the large-volume description by
Gieseker-stable pure two-dimensional sheaves, with an ample polarization
on the compact threefold $X$, and count the component supported on $S$.
This is a specified sheaf-moduli contribution, not an identification of
the full physical BPS index at the transition point.

Fix a line bundle $L$ on $S$ and let $Z\subset S$ have length $n$.
Write $i:S\hookrightarrow X$ and $\mathcal E=L\otimes I_Z$.
Every rank-one torsion-free sheaf on $S$ with determinant $L$ and second
Chern class $n$ is of this form. Its pushforward is stable for every
ample polarization: a proper subsheaf of the same rank has a nonzero
lower-dimensional quotient and hence smaller reduced Hilbert polynomial.
The moduli space on $S$ is $S^{[n]}=\operatorname{Hilb}^n(S)$,
smooth and projective of dimension $2n$ \cite{Gottsche:2003Hilbert}.

This remains a smooth component of the moduli on the compact threefold.
The normal bundle is $K_S$, and the change-of-rings sequence gives
\begin{equation}
 0\longrightarrow\operatorname{Ext}^1_S(\mathcal E,\mathcal E)
 \longrightarrow\operatorname{Ext}^1_X(i_*\mathcal E,i_*\mathcal E)
 \longrightarrow\operatorname{Hom}_S(\mathcal E,\mathcal E\otimes K_S).
 \label{eq:x9-supported-deformations}
\end{equation}
The last group vanishes: taking double duals embeds it into
$H^0(S,K_S)=0$. There are no determinant deformations, since
$H^1(S,\cO_S)=0$. Thus the first group is the $2n$-dimensional
tangent space to $S^{[n]}$. The natural map from
this proper Hilbert scheme into the stable-sheaf moduli is therefore an
isomorphism onto an open and closed smooth component. In particular,
the ambient Behrend function is $(-1)^{2n}=1$ on this component, so its
numerical Donaldson--Thomas contribution is
\begin{equation}
 \Omega_{S,L}(n)=\chi(S^{[n]})
 \label{eq:x9-supported-dt}
\end{equation}
\cite[Theorem~4.18 and Section~1.2]{Behrend:2009DT}. No other moduli
component or stability chamber is included in this definition.

The cohomology supplies a refinement
$\Omega_{S,L}(n;y)=y^{-2n}\sum_j b_j(S^{[n]})y^j$.
Here $y$ records centered cohomological degree; it is not a flavor
fugacity or an assumed infrared R-charge. We do not identify this
refinement with a protected spin character at finite gravitational
coupling: a spin-refined trace in the compact string theory need not
be invariant under changes of coupling
\cite[Section~2.6]{Manschot:2011Halos}. For these two rational
surfaces all odd Betti numbers vanish. G\"ottsche's product formula
\cite[eq.~(2.1)]{Gottsche:2003Hilbert} gives
\begin{equation}
 \begin{aligned}
 \mathcal Z_S(q,y)&=\sum_{n\geq0}\Omega_{S,L}(n;y)q^n\\
 &=\prod_{m\geq1}
 \frac{1}{(1-y^{-2}q^m)(1-q^m)^{b_2(S)}(1-y^2q^m)}.
 \end{aligned}
 \label{eq:x9-supported-refined}
\end{equation}
The answer is independent of $L$. The first numerical contributions are
\begin{equation}
 \begin{array}{c|rrrrrrr}
 n&0&1&2&3&4&5&6\\\hline
 S_-&1&4&14&40&105&252&574\\
 S_+&1&3&9&22&51&108&221
 \end{array}
 \label{eq:x9-supported-counts}
\end{equation}
Since $b_2(S_-)-b_2(S_+)=1$, even the refined ratio is
\begin{equation}
 \frac{\mathcal Z_{S_-}(q,y)}{\mathcal Z_{S_+}(q,y)}
 =\prod_{m\geq1}(1-q^m)^{-1}.
 \label{eq:x9-supported-blowup}
\end{equation}
With $q=e^{2\pi i\tau}$, the vacuum-normalized unrefined series
$q^{-\chi(S)/24}\mathcal Z_S(q,1)$ are $\eta(\tau)^{-4}$ and
$\eta(\tau)^{-3}$. Their ratio is the oscillator character
$\eta(\tau)^{-1}$ of the extra chiral boson, consistent with the wrapped-string anomaly difference
computed in \cite{Wuebben:2026quantum}. These are fixed-flux coefficients
of the rank-one string count \cite[Section~2.1]{Alim:2010M5}, not its
complete flux-summed elliptic genus.

\subsection{Shifted flux charges}
The flux shift is essential when interpreting these coefficients as
brane charges. In the convention
$\Gamma=\operatorname{ch}(i_*\mathcal E)\sqrt{\operatorname{Td}(X)}$,
Grothendieck--Riemann--Roch gives
\begin{equation}
 \begin{gathered}
 f=c_1(L)-\frac12K_S,\\
 \Gamma=\left(0,[S],i_*f,\frac12\int_S f^2+\frac{\chi(S)}{24}-n\right)
 \end{gathered}
 \label{eq:x9-supported-mukai}
\end{equation}
\cite[Appendix~B.1]{Alim:2010M5}. Thus $c_1(L)=ah+be$ on
$S_-$ gives $f=(a+\tfrac32)h+(b-\tfrac12)e$, whereas
$c_1(L)=ah$ on $S_+$ gives $f=(a+\tfrac32)h$. For the fixed representatives with restrictions
$(D_0|,D_1|,D_8|)=(0,h,-3h+e)$ on $S_-$ and
$(D_0|,D_1|,\nu|)=(0,h,-3h)$ on $S_+$, their electric
charges are
\begin{equation}
 \begin{aligned}
 (Q_{D_0},Q_{D_1},Q_{D_8})_-&=(0,a+\tfrac32,-3a-b-4),\\
 (Q_{D_0},Q_{D_1},Q_\nu)_+&=(0,a+\tfrac32,-3a-\tfrac92).
 \end{aligned}
 \label{eq:x9-supported-flux-charges}
\end{equation}
There is no integral $b$ for which the shifted exceptional flux vanishes.
Consequently \eqref{eq:x9-supported-blowup} compares oscillator counts,
not equal-charge states across the transition. A flux-summed genus
requires the shifted lattice, a quadratic refinement, and its
polarization-dependent self-dual decomposition. Nor does this sheaf
calculation determine all bound states in the physical stability
condition of the compactification.

\subsection{Magnetic pairings}
There is an additional obstruction to extracting the condensate spectrum
from this string by ordinary halo wall-crossing. For a D4--D2--D0 core
of magnetic charge $P$ and a purely electric particle
$\gamma=(0,0,\gamma_2,\gamma_0)$, choose the Dirac pairing convention
$\langle\Gamma_P,\gamma\rangle=P\cdot\gamma_2$. The compact data give
\begin{equation}
 \begin{array}{c|rrr}
 P&C_1&C_2&R=C_1+C_2\\\hline
 D_8&0&0&0\\
 D_5&0&1&1\\
 D_6&-1&0&-1
 \end{array}
 \label{eq:x9-magnetic-probes}
\end{equation}
The first row vanishes for every electric charge in
$\ZZ C_1+\ZZ C_2$, independently of the core flux and of D0 charge.
The ordinary halo index has an exponent proportional to this Dirac
pairing \cite{Denef:2007Halos}; see also
\cite[eqs.~(2.23)--(2.25)]{Manschot:2011Halos}. Its multiplier around the $D_8$
core is identically one for any spectrum supported on this charge
lattice. Equivalently, $D_5|_{S_-}=D_6|_{S_-}=0$, and every
supported flux counted above is neutral under both effective vectors.
This excludes a halo-wall-crossing determination using this core; it
does not exclude threshold bound states, contact interactions, or the
geometric transition itself. The $D_5$ and $D_6$ strings instead have
nonzero pairings with the required root and complementary nodes.
This vanishing applies to the rank-two root and nodal lattice,
not to all four finite-radius conifold charges: $D_8\iota_*e_1
=D_8\iota_*e_2=1$. The distinction will matter in
Section~\ref{subsec:x9-magnetic-higgs}.

\subsection{Moving supports and the reducible fiber}
\label{subsec:x9-moving-probes}

The nonzero pairings in \eqref{eq:x9-magnetic-probes} do not by themselves
determine a probe index. Unlike $D_8$, both $D_5$ and $D_6$ move.
Write $X=\wideX_9$ and $P_i=D_i$. Their pencils have the members
\begin{equation}
 P_5\sim D_3+D_7+D_8,\qquad P_6\sim D_4+D_7+D_8,
 \qquad P_5^2=0,\quad P_6^2=C_1.
 \label{eq:x9-probe-pencils}
\end{equation}
The first defines the K3 fibration.
A smooth member of the second is a K3 surface blown up at a point,
with canonical curve $C_1$. Each pencil has two sections, represented
in Cox coordinates by $(z_5,z_3z_7z_8)$ and $(z_6,z_4z_7z_8)$.
The second has scheme-theoretic base locus $C_1$, a smooth complete
intersection with normal bundle $\cO(-1)\oplus\cO(-1)$.
Its universal member is therefore $\operatorname{Bl}_{C_1}X$.
Counting a single fixed K3 would omit the motion of the support
\cite[Section~2.2]{Bouchard:2017Vertical}.

The higher cohomology of $\cO_X(P_i)$ vanishes and
$h^0(\cO_X(P_i))=2$. The divisor sequence then gives
$H^1(T,\cO_T)=0$ for every member. The structure-sheaf seed
charges, whether or not their full pencil is a stable moduli
component, are
\begin{equation}
 \Gamma_{5,0}=(0,P_5,0,1),\qquad
 \Gamma_{6,0}=(0,P_6,-\tfrac12C_1,\tfrac{11}{12}).
 \label{eq:x9-moving-seed-charges}
\end{equation}
Appendix~\ref{app:pencil} computes zero- and one-point pencil
contributions under a separate reduced-member hypothesis.
The support restrictions and isolated contributions below are
unconditional for the specified resolved geometry.

\paragraph{Where the missing probe states must occur.}
For $P_5$, the limitation is stronger than motion of the support.
Any effective support cycle of class $P_5$ is vertical for $\rho$:
nefness and $P_5^2=0$ imply that each component has zero intersection
with $P_5$ and an ample divisor. A sheaf on such a support is unchanged
by tensoring with $\cO_X(P_5)$, which is pulled back from the base and
trivial over a neighborhood of the finitely many supporting fibers.
Its induced D2 charge therefore obeys
\begin{equation}
 Q_{P_5}=P_5\cdot\operatorname{ch}_2(\mathcal F)=0.
 \label{eq:x9-vertical-charge-selection}
\end{equation}
Adding $C_2$ or $R$ to a seed changes this charge by one. Such a
core--particle state cannot lie in any pure two-dimensional sheaf
component of magnetic charge $P_5$ in this large-volume description.
It would require a different physical stability description, for
example a derived object; this is not a no-go for the physical state.

For $P_6$, set $A_6=D_7-D_6=-D_4-D_8$. Every pencil member except
$T_*=D_4+D_7+D_8$ is disjoint from $D_4$ and $D_8$.
Thus $\cO_X(A_6)$ restricts trivially to it, and all sheaves supported
on those members have $Q_{A_6}=0$, including the seeds. But
\begin{equation}
 A_6\cdot C_1=0,\qquad A_6\cdot C_2=A_6\cdot R=1.
 \label{eq:x9-reducible-probe-selection}
\end{equation}
Within this pencil, a sheaf obtained by adding $C_2$ or $R$ must
therefore be supported on $T_*$. This reduces the sheaf calculation
to a reducible surface, with gluing along its double curves. The vertical modular formulas of
\cite{Bouchard:2017Vertical} assume irreducible fibers and cannot be
applied directly to this contribution. Neither a fixed smooth-K3
oscillator series nor the seed components above determine it.

\paragraph{An isolated contribution with cooperative electric charge.}
Isolated components of this reducible-surface moduli problem can be computed
without the additional pencil hypothesis. Put $T=T_*$ and
$\mathcal I=I_{C_2,T}$. Consider the two sheaves
\begin{equation}
 \mathcal F_2=\mathcal Hom_T(\mathcal I,\cO_T),\qquad
 \mathcal F_R=\mathcal Hom_T(\mathcal I,\omega_T)
             =\mathcal F_2\otimes\cO_X(P_6).
 \label{eq:x9-glued-sheaves}
\end{equation}
Sheaves on $T$ are again understood to be pushed forward to $X$.
Both are $J_*$-Gieseker stable and define isolated reduced points of
the stable-sheaf moduli on $X$. Their charges and numerical DT
contributions are
\begin{equation}
 \begin{array}{c|c|c}
 &\Gamma&\Omega^{\rm iso}\\\hline
 \mathcal F_2&(0,P_6,-\tfrac12 C_1+C_2,-\tfrac1{12})&1\\
 \mathcal F_R&(0,P_6,\tfrac12 C_1+C_2,-\tfrac1{12})&1 .
 \end{array}
 \label{eq:x9-glued-indices}
\end{equation}
In particular,
$\Gamma(\mathcal F_R)-\Gamma_{6,0}=(0,0,R,-1)$.
The superscript denotes these individual components, not the full
invariants at their charges.

The construction also determines the gluing at the double curve.
The complete intersection $C_2=D_2D_4\simeq\mathbb P^1$ has
normal bundle $\cO(-1)\oplus\cO(-1)$ in $X$, and
\begin{equation}
 (D_4,D_7,D_8)\cdot C_2=(-1,1,0),\qquad P_6\cdot C_2=0.
 \label{eq:x9-gluing-incidence}
\end{equation}
It meets $D_4\cap D_7$ at one transverse point $p$, away from $D_8$.
Near $p$ the surface and ideal have the form
$\cO_T=\CC[[x,y,z]]/(xy)$ and $\mathcal I=(x,z)$.
Thus $C_2$ is not Cartier on $T$: the dual ideal in
\eqref{eq:x9-glued-sheaves} is rank-one torsion-free, but not locally
free at $p$. This is essential. The naive restrictions
$(\cO_{D_4}(C_2),\cO_{D_7},\cO_{D_8})$ have degrees one and zero
on the two sides of $D_4\cap D_7$, and cannot be glued as a line
bundle.

The ideal $\mathcal I$ is maximal Cohen--Macaulay on the Gorenstein
surface $T$, by the sequence
$0\to\mathcal I\to\cO_T\to\cO_{C_2}\to0$ and the depth lemma.
Duality gives pure sheaves in \eqref{eq:x9-glued-sheaves} and sequences
\begin{equation}
 \begin{aligned}
 0&\longrightarrow\cO_T\longrightarrow\mathcal F_2
       \longrightarrow\cO_{C_2}(-2)\longrightarrow0,\\
 0&\longrightarrow\omega_T\longrightarrow\mathcal F_R
       \longrightarrow\cO_{C_2}(-2)\longrightarrow0.
 \end{aligned}
 \label{eq:x9-glued-extensions}
\end{equation}
Here $\omega_T=\cO_T(P_6)=\cO_T(C_1)$, and its restriction to
$C_2$ is trivial. Since $\chi(\cO_{C_2}(-2))=-1$, these sequences
give precisely the D0 shift in \eqref{eq:x9-glued-indices}; it must
not be discarded when interpreting the electric flux as an added
particle.

For stability, a proper subunion $U\subset T$ has the same generic
double-curve gluing as $\cO_T$. For $T=U+(T-U)$ define
\[
 \delta_T(U)=[J_*U(2T-U)](J_*^2T)-(J_*T^2)(J_*^2U).
\]
Relative to $\cO_T$, the component
fluxes of $\mathcal F_2$ are $(C_2,0,0)$ and those of
$\mathcal F_R$ are $(C_2,C_1,0)$, in the order $(D_4,D_7,D_8)$.
Let $\gamma$ be the total flux and $\gamma_U$ its sum over $U$.
The numerator comparing the maximal subsheaf supported on $U$ with
the full sheaf is
\begin{equation}
 \delta_{\gamma}(U)=\delta_T(U)
   +2\big[(J_*\cdot\gamma)(J_*^2\cdot U)
          -(J_*\cdot\gamma_U)(J_*^2\cdot T)\big].
 \label{eq:x9-glued-stability}
\end{equation}
Point defects do not change this coefficient. For the six subunions
$1,2,3,12,13,23$, the numerators are
\begin{equation}
 \begin{array}{c|rrrrrr}
 \mathcal F_2&3388&5222&1488&1482&4282&4334\\
 \mathcal F_R&3630&4974&1494&1476&4530&4092 .
 \end{array}
 \label{eq:x9-glued-stability-data}
\end{equation}
All are strictly positive. Full-multirank proper subsheaves have a
nonzero lower-dimensional quotient, hence a smaller reduced Hilbert
polynomial. This proves stability, including the non-locally-free
gluing point.

To compute the ambient deformation space, let $I_{C_2}$ be the ideal
in $X$. Since $C_2$ is disjoint from the pencil base curve $C_1$,
the restriction of the two sections of $\cO_X(P_6)$ to
$\cO_{C_2}(P_6)=\cO_{C_2}$ is surjective. Together with the higher
cohomology vanishing above, this gives
$R\operatorname{Hom}(\cO_X(-P_6),I_{C_2})=\CC$ in degree zero.
Consequently
\begin{equation}
 0\longrightarrow\cO_X(-P_6)\longrightarrow I_{C_2}
       \longrightarrow\mathcal I\longrightarrow0
 \label{eq:x9-glued-twist}
\end{equation}
is the spherical-twist triangle. The autoequivalence of
\cite[Theorem~1.2]{Seidel:2001Twists} identifies the self-Ext groups
of $I_{C_2}$ and $\mathcal I$. The former has no first-order
deformations: $H^0(N_{C_2/X})=0$ and $H^1(\cO_X)=0$.
Derived duality satisfies
$R\mathcal Hom_X(\mathcal I,\cO_X)=\mathcal F_R[-1]$;
it, and tensoring by $\cO_X(-P_6)$, preserve the dimension of
self-$\operatorname{Ext}^1$. Thus
\begin{equation}
 \operatorname{Ext}^1_X(\mathcal F_2,\mathcal F_2)
 =\operatorname{Ext}^1_X(\mathcal F_R,\mathcal F_R)=0.
 \label{eq:x9-glued-rigidity}
\end{equation}
Each stable point is therefore an open and closed smooth
zero-dimensional component, with Behrend weight $+1$
\cite[Section~1.2 and Theorem~4.18]{Behrend:2009DT}.

The second sheaf supplies a compact contribution carrying the
cooperative electric class. It remains distinct from the desired five-dimensional particle
index. The complete moduli at either charge have not been classified,
and physical stability at the cooperative contraction is not established.
Section~\ref{subsec:x9-probe-transport} extends the Gieseker-stability
calculation and explains why the raw D0 shift is frame-dependent.
At these fixed charges, the divisor
sequence and Serre duality give
\[
 R\operatorname{Hom}_X(\cO_{C_2}(-1),\cO_T)=0.
\]
Thus replacing the quotient in
\eqref{eq:x9-glued-extensions} by the zero-D0 curve sheaf does not
give a nontrivial extension of $\cO_T$. This excludes that elementary
construction, not other sheaves or derived objects at zero D0 shift.
